\section{Rekwizyty}

\subsection{Kostki}
\textit{Labirynt śmierci} składa się z niniejszej instrukcji oraz 346 żetonów.
Niezbędne jest również zaopatrzenie się w jedną lub kilka zwykłych, sześciennych kostek do gry.
Gdyby było to niemożliwe, gracze mogą korzystać z żetonów oznaczonych numerami od 1 do 6, umieszczając je w pojemniku (np. w plastikowym kubeczku) i losowo wyciągając jeden z nich zawsze wtedy, gdy potrzebny jest rzut kostką.
Jeżeli przepisy gry wymagają rzutu dwoma lub trzema kostkami (takie sytuacje się zdarzają), można po prostu ciągnąć żeton dwu- lub trzykrotnie.
Należy jednak przy tym pamiętać, aby po każdym wyciągnięciu żeton wrócił z powrotem do pojemnika.

\subsection{Plansza}
Gra nie posiada planszy. Gracze sami ją konstruują, ustawiając jeden za drugim odpowiednie żetony w miarę eksploracji labiryntu.
Ustawiane kolejno żetony muszą wzajemnie do siebie pasować tzn. korytarz musi łączyć się z korytarzem, a przejście z przejściem.
W dalszej części instrukcji jeden taki żeton (czyli komnata lub fragment korytarza) nazywany będzie segmentem labiryntu.

\subsection{Żetony}
346 załączonych do gry żetonów reprezentuje korytarze, komnaty (puste lub zawierające znaleziska), pastacie, które tworzą oddział oraz potwory.
Dołączonych jest również 6 żetonów zastępujących kostkę. Przykłady:

TUTAJ BĘDZIE OBRAZEK Z PRZYKŁADAMI ŻETONÓW

\subsection{Modyfikacje zasad}
Zasady gry są tak skonstruowane, by umożliwić graczon ich dowolną modyfikację lub formułowanie własnych, nie występujących w instrukcji przepisów.
Gracze nie muszą sztywno trzymać się tych zasad, które zostały im przedstawione.
Wręcz przeciwnie~--- zachęcamy do ich zmieniania, jeżeli wszyscy grający uznają, że taka zmiana przyniesie korzyść grze.

Ta dość duża swoboda interpretacji przepisów jest spowodowana m.in. tym, że \textit{Labirynt śmierci} należy do tego typu gier, w których niemożliwością jest sztywne ujęcie zasadami wszystkich sytuacji, jakie mogą się pojawić w toku rozgrywki. W sytuacjach niejasnych i nie wyjaśnionych w instrukcji należy kierować się po prostu zdrowym rozsądkiem lub (jeśli ktoś woli) rzutem kostką.
